\documentclass[10pt,letterpaper]{article}
\usepackage{cornell-notes}

\title{Clause 05.02.10: View Methods}
\author{}
\date{}
\setDocTitle{Clause 05.02.10: View Methods}
\setDocSubtitle{ISO/IEC/IEEE 42010:2022}
\setDocOwner{ISO 42010 Cornell Notes}
\setDocDate{\today}
\setCornellCollection{ISO/IEC/IEEE 42010:2022 Cornell Notes}
\setCornellUnitType{Clause}
\setCornellUnitNumber{05.02.10}
\setCornellUnitTitle{View Methods}
\hypersetup{pdftitle={Clause 05.02.10: View Methods},pdfsubject={Cornell notes},pdfauthor={}}

\begin{document}
\maketitle
\begin{CornellOverviewBox}{Overview}
{\Large\bfseries\textcolor{CNavy}{Section 5.2.10: View methods}}\\[3pt]
{\small\textbf{Classification:} Conceptual foundation}\\
{\small\textbf{Learning objective:} Explain the role of view methods and connect it to related architecture-description concepts.}
\end{CornellOverviewBox}
\vspace{3pt}

\begin{CornellNotesTable}
\CornellNoteRow{Purpose / key concept}{A specification of an architecture viewpoint includes one or more view methods.}
\CornellNoteRow{Core details}{\begin{itemize}[leftmargin=*,nosep,topsep=1pt]
\item View methods provide guidance, heuristics, metrics, patterns, design rules or guidelines, best practices and examples to aid in view construction and use of associated views.
\item View methods specify expression rules, modelling methods, analysis techniques and other operations on views.
\item These methods specify the AD elements used when creating the view and methods to analyse, interrogate or query views to assess properties of interest.
\item Requirements on view methods are specified in 8.3.
\end{itemize}}
\CornellNoteRow{Example / application}{View methods pertaining to: chaining dependencies to assess the impact of a change; workshops to trade-off qualities or other concerns; boundary analyses to determine whether context and entity of interest are well defined; guidance on partitioning; analysis of architectural complexity; creation and enforcement of architecture styles (such as layered, aspect-oriented); pattern families to promote intended properties of the entity of interest; analyses against requirements for completeness and coverage; ...}
\CornellNoteRow{Interpretive note}{\begin{itemize}[leftmargin=*,nosep,topsep=1pt]
\item View methods are usually defined in a viewpoint and are referenced by or used by model kinds, architecture description frameworks, and architecture description languages.
\end{itemize}}
\CornellNoteRow{Related sections}{8.3}
\CornellNoteRow{Active recall}{\begin{itemize}[leftmargin=*,nosep,topsep=1pt]
\item What is the central role of View methods?
\item How does View methods connect to adjacent architecture-description concepts?
\end{itemize}}
\end{CornellNotesTable}


\begin{CornellSummaryBox}{Summary}
A specification of an architecture viewpoint includes one or more view methods.
\end{CornellSummaryBox}

\end{document}
